\documentclass[11pt,letterpaper]{article}

% Standalone source: upload this file to Overleaf and compile with pdfLaTeX.
\usepackage[T1]{fontenc}
\usepackage[utf8]{inputenc}
\usepackage{charter}
\usepackage[margin=1in]{geometry}
\usepackage{microtype}
\usepackage[hidelinks]{hyperref}
\hypersetup{
  pdftitle={Research Statement},
  pdfauthor={Mst Afsana Mimi},
  pdfsubject={PhD application, University of Miami, Fall 2027}
}
\input{glyphtounicode}
\pdfgentounicode=1
\pagestyle{empty}
\setlength{\parindent}{0pt}
\setlength{\parskip}{0.75\baselineskip}
\linespread{1.04}
\frenchspacing

\begin{document}

\begin{center}
  {\Large\bfseries Research Statement}\\[6pt]
  {\large Mst Afsana Mimi}\\[4pt]
  {\small Prospective PhD Student, Fall 2027}
\end{center}
\vspace{2pt}

For my PhD, I would like to explore data-efficient predictive modeling and uncertainty quantification for complex dynamical systems. Building on my experience with Gaussian process regression, I am interested in how probabilistic models can learn from limited data and express uncertainty in their predictions. I would also like to investigate how generative models and simulation can be used together to improve predictions when collecting real-world data is difficult or expensive. My current work on agentic AI for relational prediction has also made me curious about how agents can select useful information and decide what to do next. In particular, I want to understand how uncertainty can guide the choice of new observations or simulations. I am also interested in evaluating how well these models perform when the available data do not fully represent the conditions they encounter. I hope to explore how these ideas can help learning systems make reliable predictions and informed decisions with limited data.

\end{document}
